\section{The Causal Inference Problem}

The quest for rigorous causal reasoning traces back to the pioneering work of John Snow during the 1854 London cholera outbreak~\cite{snow1855mode}, but its formalization did not happen until in the 20th century, when Jerzy Neyman introduced the potential outcomes notation in 1923 for agricultural experiments~\cite{neyman1923application}, Ronald Fisher developed the theory of randomized controlled trials~\cite{fisher1971design}, and Donald Rubin extended the potential outcomes framework to observational studies decades later~\cite{rubin1974estimating}.
Armed with these rigorous frameworks, causal inference has transformed decision-making across virtually every domain of human inquiry.
In medicine, randomized controlled trials became the gold standard for drug approval, ushering in the era of evidence-based medicine that has saved countless lives~\cite{sackett1996evidence}.
In economics, causal methods revolutionized policy evaluation: Card and Krueger's difference-in-differences study on minimum wage~\cite{card1993minimum} overturned decades of conventional wisdom, while Angrist and Imbens formalized instrumental variable methods for extracting causal effects from observational data~\cite{angrist1996identification}.
The field's crowning recognition came in 2021, when the Nobel Prize in Economic Sciences was awarded jointly to David Card, Joshua Angrist, and Guido Imbens \emph{``for their methodological contributions to the analysis of causal relationships''}---a testament to how profoundly these methods have reshaped empirical research~\cite{nobel2021economics}.
Beyond academia, tech giants now run millions of A/B tests annually to causally evaluate product changes~\cite{kohavi2020trustworthy}, while public health agencies relied on causal inference frameworks to rapidly evaluate COVID-19 vaccine efficacy during the pandemic~\cite{polack2020safety}.
With research ranging from policy-making to medicine to application user experience, the ability to distinguish causation from correlation has become indispensable to modern science and industry alike.

Conceptually, causation can be defined through \emph{counterfactuals}: $X$ causes $Y$ if $Y$ would have been different had $X$ not occurred.
Yet, we can never observe both realities for the same unit---a patient either receives a drug or does not; a developer either adopts a tool or does not; a project uses one programming language or another.
This ``fundamental problem of causal inference,'' that individual causal effects are inherently unobservable, means that causal claims can never be directly verified.
Different intellectual traditions have developed distinct approaches to navigate this challenge.
Specifically, three conceptual frameworks have shaped modern causal inference, each offering a distinct lens for reasoning about cause and effect:
\begin{itemize}[leftmargin=15pt]
    \item \textbf{The Potential Outcomes Framework:} Pioneered by Neyman in 1923 and formalized by Rubin~\cite{rubin1974estimating}, this framework defines causation through \emph{counterfactual comparisons}. 
    For each unit $i$, there exist two potential outcomes: $Y_i(1)$ (outcome if treated) and $Y_i(0)$ (outcome if not treated). 
    The individual causal effect is $Y_i(1) - Y_i(0)$, but since only one outcome is ever observed, the framework focuses on estimating population-level quantities like the Average Treatment Effect $ATE = E[Y(1) - Y(0)]$ or the Average Treatment Effect on Treated $ATT =E[Y_i(1)-Y_i(0)|i \text{ treated}]$.
    The central question becomes: Under what conditions can we use observed data to estimate these counterfactual quantities?
    Within this framework, randomized controlled trials are the gold standard, as asymptotic theory shows that randomized assignment ensures unbiased estimation of $ATE$ by design.\footnote{The actual proof would need a few additional assumptions we will not cover in this thesis, see~\citet{cunningham2021causal}.}
    Many statistical modeling approaches can also be examined through the lens of potential outcomes, some with assumptions that can be explicitly tested and others with assumptions that are untestable but potentially plausible.
    
    \item \textbf{The Graphic Causal Models:} Developed by Pearl~\cite{pearl2009causality}, this framework represents causal relationships as directed acyclic graphs (DAGs), where nodes are variables and edges represent direct causal effects.
    Its strength lies in making causal assumptions visually explicit: A DAG encodes what the researcher believes about the data-generating process, thereby allowing the researcher to directly incorporate alternative explanations and common threats to validity into causal inference designs.
    For example, when investigating the causal relationship between $X$ and $Y$, a \emph{confounder} $Z$ (one that affects both $X$ and $Y$) can be directly encoded as $X \leftarrow Z \rightarrow Y$.
    With an explicitly defined causal DAG, researchers can assess the plausibility of their assumptions even when a direct statistical test is unavailable; critics can challenge the validity of findings by pointing to specific edges, missing nodes, missing paths, or alternative structures, forming a constructive ground for debates and improved designs in future research.\footnote{What we present here is a useful oversimplification of the full structural causal model framework, which many applied scientists use in combination with potential outcomes and design-based arguments to justify their research designs. See~\citet{pearl2009causality} for a more rigorous coverage and formalization of structural causal modeling approaches.}
    
    \item \textbf{The Design-Based (Credibility Revolution) Approach:} Championed by applied economists like Angrist and Imbens~\cite{angrist1996identification}, this approach is less a formal framework and more a philosophy of pragmatic empirical practice.
    Rather than starting from abstract potential outcomes or graphical models, it asks: What feature of the real world generates plausible \emph{natural experiments} with exogenous variation in treatment?
    The emphasis is on identifying situations in which nature, policy, or institutional rules create variation that approximates random assignment and on transparently arguing why alternative explanations are implausible.
    A study's credibility rests not on mathematical formalism but on whether the identification strategy would convince a skeptical reader that confounding has been adequately addressed.
\end{itemize}
While these frameworks differ in formalism and emphasis, they share a common commitment: \emph{Making causal assumptions explicit and assessable so they can be examined and debated, rather than left implicit and buried in the underlying statistical models.}
In fact, these frameworks are often used together in applied research as a basis for \emph{causal credibility assessments}.
Critics can then interrogate whether that assumption is plausible in the specific research context, rather than debating vague notions of ``correlation versus causation'' or whether the findings ``seem valid.''
Depending on their assumptions, these frameworks yield a family of \emph{identification methods}: research designs that, under their stated assumptions, provide a complete path from observed data to causal effects.
Alongside these identification methods, researchers also frequently employ \emph{control methods}---techniques like fixed effects and covariate adjustment that address specific sources of confounding without claiming complete identification.

To illustrate the promises and challenges of causal reasoning in software engineering, consider a deceptively simple question: Does the choice of programming language causally affect defect-proneness?
A large-scale study~\cite{DBLP:conf/sigsoft/RayPFD14, DBLP:journals/cacm/RayPDF17} analyzed millions of commits across thousands of GitHub projects and reported that functional and statically-typed languages are associated with fewer defects.
A subsequent reproduction study~\cite{DBLP:journals/toplas/BergerHMVV19}, however, found that effect sizes shrank dramatically after correcting methodological issues, sparking a series of rebuttals~\cite{DBLP:journals/corr/abs-1911-07393, DBLP:journals/corr/abs-1911-11894}.
The core difficulty lies in the nature of the data: Projects written in different languages differ systematically in ways that correlate with defect rates---team expertise, project maturity, application domain, and coding conventions all vary with language choice.
Observing that Haskell projects have fewer bugs than C projects tells us little about what would happen if the same team rewrote the same application in a different language.
This reflects a fundamental challenge in causal inference, \emph{confounding}: Observed associations may reflect selection effects rather than causal mechanisms.
To properly address confounding, a researcher can either: (1) design a study that provides full causal identification by design, such as randomized controlled trials proven through the potential outcome framework; (2) run a large-scale observational study and use causal DAGs to assess the plausibility of causal interpretation from their findings; (3) find natural experiments where it is plausible to argue that the treatment is as-good-as-random in the specific setting.
We will discuss in Chapter~\ref{chp:proposed-work} regarding how to address confounding in this specific problem.

Importantly, in many research settings, and perhaps more so in software engineering, there is an inherent trade-off between internal and external validity across research designs.
At one end of the spectrum, randomized experiments, while eliminating all sources of confounding, may be impractical, unethical, or---crucially---may alter the setting so fundamentally that findings fail to generalize.
For example, a controlled lab experiment comparing how developers implement the same task in different programming languages may achieve perfect internal validity but tell us little about the actual effect in the messy reality of production environments with deadlines, legacy code, and team dynamics.
At the other end, observational studies of real-world behavior preserve ecological validity but are much more challenging to rule out confounders and alternative explanations.
Between these extremes lies a hierarchy of quasi-experimental methods that attempt to find or construct a setting comparable to an actual experiment, but the attempt to construct a quasi-experimental setting alone may limit the extent to which the findings can generalize to other settings.  
No single study, however well-designed, can definitively establish a causal claim across all contexts, and confidence in a causal relationship grows through \emph{triangulation}---when multiple studies using different designs, data sources, and populations converge on consistent findings.
The credibility of a causal theory ultimately rests not on any one method but on the accumulated weight of evidence across complementary approaches.

The remainder of this chapter develops the technical foundations for causal inference in software engineering research.
Section~\ref{sec:pitfalls} examines common pitfalls in descriptive and correlational research, explaining why correlation does not imply causation and why multivariate regression alone cannot establish causal claims in most settings.
Section~\ref{sec:formalism} formalizes the key concepts, presenting the potential outcomes framework and graphical causal models that are tremendously useful in reasoning about causality and assessing the credibility of causal claims.
%Section~\ref{sec:identification} surveys identification strategies---randomized controlled trials, difference-in-differences, regression discontinuity, and instrumental variables---that provide complete causal identification under stated assumptions.
%Section~\ref{sec:control} discusses control strategies---matching, fixed effects, and generalized method of moments---that address specific sources of confounding without providing complete identification.
Section~\ref{sec:threats} discusses threats to validity in causal designs and the role of sensitivity analysis in strengthening credibility.
Finally, Section~\ref{sec:pragmatic} advocates a pragmatic approach to causal inference and how to conduct explicit causal credibility assessments in empirical software engineering research.

\section{Pitfalls in Descriptive and Correlational Research}
\label{sec:pitfalls}

Before presenting the causal inference approach, it is instructive to examine why naive approaches often fail to support causal claims.
Understanding these pitfalls not only clarifies the challenges researchers must overcome to make a causal claim but also reveals why the solutions are non-trivial---why simply ``being careful'' or ``controlling for more variables'' cannot substitute for genuine causal reasoning and explicit assessments of causal credibility.

\subsection{Why Correlation Does Not Imply Causation}

The dictum ``correlation does not imply causation'' is perhaps the most frequently cited methodological principle in empirical research, yet it remains the most frequently violated.
The violation persists not because researchers are careless, but because the underlying mechanisms that generate spurious correlations are subtle and pervasive.
Let us examine these mechanisms through the lens of our programming language example.

Recall the question: Does choice of programming language causally affect defect-proneness?
Suppose we compute the correlation between language type (say, functional vs. imperative) and defect density across thousands of GitHub projects and find that functional language projects have 20\% fewer defects.
Can we conclude that switching to a functional language would reduce defects by 20\%?
The answer is almost certainly no, and understanding why requires distinguishing three mechanisms that can generate spurious correlations.

\paragraph{Confounding: The Hidden Common Cause.}
Confounding occurs when a third variable $Z$ causally affects both the treatment $X$ and the outcome $Y$, inducing an association between them even when no direct causal relationship exists.
In our programming language example, consider \emph{developer expertise}.
It is very likely that experienced developers are both more likely to choose functional languages (perhaps because these require more sophisticated programming knowledge) and more likely to write low-defect code (because expertise directly reduces errors).
If the above mechanism were true, the observed correlation between functional languages and fewer defects may simply reflect this selection effect: We are comparing code written by different types of developers, not the causal effect of the language itself.

To see this mathematically, let $Y$ denote defect density, $X$ denote language choice (1 for functional, 0 for imperative), and $Z$ denote developer expertise.
The true data-generating process might be:
\begin{align}
    X &= f(Z, U_X) \quad \text{(language choice depends on expertise)} \\
    Y &= g(Z, U_Y) \quad \text{(defects depend on expertise, not language)}
\end{align}
where $U_X$ and $U_Y$ are independent error terms.
Even though $X$ has no causal effect on $Y$, the correlation $\text{Cor}(X, Y) \neq 0$ because both are driven by the common cause $Z$.
More precisely, decomposing the observed association:
\begin{equation}
    E[Y | X=1] - E[Y | X=0] = \underbrace{\tau}_{\text{causal effect}} + \underbrace{(E[Z | X=1] - E[Z | X=0])}_{\text{selection difference}} \cdot \underbrace{\frac{\partial Y}{\partial Z}}_{\text{effect of $Z$}}
\end{equation}
If the causal effect $\tau = 0$ but functional language adopters have higher expertise ($E[Z|X=1] > E[Z|X=0]$) and expertise reduces defects ($\partial Y / \partial Z < 0$), we observe a negative correlation between functional languages and defects---but it is a pure artifact of selection.

\paragraph{Reverse Causality: Getting the Arrow Backwards.}
Reverse causality occurs when the observed association reflects $Y$ causing $X$ rather than $X$ causing $Y$.
In our example, perhaps teams whose projects have persistent quality problems switch to more ``rigorous'' languages like Haskell in an attempt to impose discipline.
If this were true, we might observe that Haskell projects have \emph{more} defects---not because Haskell causes defects, but because defect-prone teams select into Haskell.
Alternatively, successful low-defect projects may attract functional programming enthusiasts who refactor the codebase, creating the opposite spurious correlation.

The fundamental problem is that cross-sectional data cannot distinguish between $X \to Y$ and $Y \to X$.
Both causal structures generate the same correlation.
Only temporal information (did language choice precede defect reduction?) combined with careful reasoning about the selection process can help disentangle these possibilities.

\paragraph{Selection Bias: Conditioning on a Consequence.}
Selection bias arises when the sample is not representative of the population of interest, or when the data-generating process induces spurious associations through conditioning on a \emph{collider}---a variable caused by both $X$ and $Y$.

Consider an extreme example: Suppose we analyze only ``successful'' GitHub projects (those with more than 1,000 stars).
A project becomes successful if it either (a) solves a valuable problem despite bugs, or (b) has exceptional code quality.
Among successful projects, we might observe a \emph{negative} correlation between functional languages and defects: Successful functional projects may have achieved success through quality, while successful imperative projects may have succeeded despite bugs by solving critical problems.
This correlation is entirely an artifact of conditioning on success---it tells us nothing about the causal effect of language choice.

More subtly, analyzing only projects that ``survived'' to a certain age conditions on a post-treatment variable (survival), potentially inducing bias.
If functional language projects fail faster when they have bugs (because functional programming communities have higher quality standards), the surviving functional projects are a selected sample, and comparisons with surviving imperative projects are confounded.

\subsection{Why Multivariate Regression (Usually) Does Not Suffice}

Confronted with confounding, researchers often adopt what seems like an obvious solution: ``control for'' the confounders using multivariate regression.
If developer expertise confounds the relationship between language and defects, why not simply include expertise as a covariate?
While intuitive, this approach has fundamental limitations that render it insufficient for causal inference.

\paragraph{Omitted Variable Bias: The Inescapable Problem of Unobserved Confounders.}
The most fundamental limitation is that regression can only adjust for \emph{observed} confounders.
In practice, important confounders are often unmeasured or unmeasurable.
Returning to our example, consider the confounders that plausibly affect both language choice and defect rates: Developer expertise, yes, but also team culture, project complexity, application domain, code review practices, testing discipline, deadline pressure, and organizational incentives.
Can we measure all of these from repository data in a convincing way?
Probably not.

The mathematical consequence is stark.
Suppose the true model is:
\begin{equation}
    Y_i = \alpha + \tau X_i + \beta Z_i + \gamma U_i + \epsilon_i
\end{equation}
where $Z$ is an observed confounder and $U$ is an unobserved confounder.
If we regress $Y$ on $X$ and $Z$ only, the estimated coefficient on $X$ is:
\begin{equation}
    \hat{\tau} \xrightarrow{p} \tau + \gamma \cdot \frac{\text{Cov}(X, U | Z)}{\text{Var}(X | Z)}
\end{equation}
The bias term $\gamma \cdot \text{Cov}(X, U | Z) / \text{Var}(X | Z)$ depends on the relationship between treatment and the omitted confounder \emph{after} controlling for $Z$.
Unless $\text{Cov}(X, U | Z) = 0$---which would require $Z$ to fully mediate the relationship between $X$ and $U$---the estimate remains biased.
Adding more controls helps only if those controls are themselves confounders; adding irrelevant controls simply reduces precision without reducing bias.

%The uncomfortable implication is that \emph{no amount of control variables can substitute for a credible identification strategy}.
%A regression with 50 covariates is no more causally credible than one with 5 if the key confounders remain unobserved.
%The number of controls is irrelevant; what matters is whether the remaining variation in treatment is as-if random.

\paragraph{Bad Controls: When Adjustment Makes Things Worse.}
A subtler problem is that not all variables \emph{should} be controlled for, even if they are observed.
Including the wrong controls can introduce bias where none existed, or amplify existing bias.
Two cases are particularly important:

\begin{itemize}[leftmargin=15pt]
    \item \emph{Mediator bias.}
    A mediator is a variable on the causal pathway from treatment to outcome ($X \to M \to Y$).
    Controlling for a mediator blocks part of the causal effect and biases the estimate toward zero.
    In our example, suppose functional languages reduce defects partly by encouraging immutability, which in turn prevents certain bug classes.
    If we ``control for'' use of immutable data structures, we block this pathway and underestimate the true effect of language choice on defect proneness.
    \item \emph{Collider bias.}
    A collider is a variable caused by both treatment and outcome ($X \to C \leftarrow Y$).
    Conditioning on a collider \emph{creates} spurious associations between its causes, even when they are independent.
    Suppose a project's ``popularity'' depends on both the novelty of its language choice and its quality.
    Controlling for popularity induces a negative correlation between language novelty and quality among popular projects: If we know a project is popular, learning that it uses a novel language leads us to infer it must have lower quality (otherwise, why would quality alone not suffice for popularity?).
    This is a pure statistical artifact.
\end{itemize}
The lesson is that without a causal model specifying the data-generating process, researchers cannot determine which controls are appropriate.
A regression with 100 covariates and 0.99 $R^2$ is not better than others in terms of causal credibility.
Regression software will happily estimate coefficients regardless of whether the specification is causally meaningful. 

\begin{comment}
\paragraph{Functional Form Misspecification.}
An even subtler problem for multivariate regressions is that most standard regressions assume linear, additive relationships.
If the true effect is nonlinear, if treatment effects are heterogeneous, or if confounders interact with treatment, the regression coefficient may not correspond to any meaningful causal quantity.
In our example, the effect of language choice may depend on project size, team composition, or application domain.
A single coefficient averaging across all contexts may be misleading for any particular context.

\paragraph{The Table 2 Fallacy: Misinterpreting ``Controlled'' Estimates.}
These problems manifest in a pervasive interpretive error that Westreich and Greenland term the ``Table 2 Fallacy''~\cite{westreich2013table}: The belief that if a coefficient ``remains significant after controlling for $Z$,'' the association must be causal.
This interpretation fails for multiple reasons:
\begin{itemize}[leftmargin=15pt]
    \item If $Z$ is a mediator, the controlled estimate underestimates the true effect.
    \item If $Z$ is a collider, conditioning on it introduces spurious association.
    \item If $Z$ is a confounder but imperfectly measured, residual confounding persists.
    \item If important confounders remain omitted, the estimate is still biased regardless of what was included.
\end{itemize}
The interpretation of ``controlling for $Z$'' depends entirely on the causal structure relating $X$, $Y$, and $Z$---a structure that regression alone cannot reveal.
\end{comment}

\subsection{When Can Multivariate Regression Support Causal Claims}

The preceding discussion may leave the impression that multivariate regression is never useful for causal inference.
This would be an overstatement---and an impractical one.
The formal condition under which regression identifies causal effects is \emph{selection on observables} (also called conditional ignorability): $(Y(1), Y(0)) \perp D \mid X$, meaning that conditional on observed covariates, treatment is as-if randomly assigned.
This requires that all confounders be measured and correctly included, that no post-treatment variables contaminate the model, and that there is sufficient overlap between the treated and control groups.
Methods such as propensity score matching, inverse probability weighting, and regression adjustment all rest on the same assumption — only their implementation differs, not the fundamental identification claim.

\paragraph{The Role of Theory.}
This is where causal inference rejoins substantive domain knowledge.
If a researcher can articulate a complete causal model---a DAG representing all relevant variables and their relationships---and the regression specification correctly implements the adjustment set dictated by that DAG, then the coefficient \emph{does} have a causal interpretation.
The regression is not merely ``controlling for stuff''; it operationalizes a theoretical claim about the data-generating process.
The credibility of the causal estimate then rests entirely on the credibility of the underlying theory: Have we correctly identified all common causes of treatment and outcome?
Have we avoided conditioning on mediators and colliders?
Is the proposed causal structure consistent with domain knowledge and prior evidence?

In this sense, \emph{all} causal inference from observational data is theory-dependent.
Design-based methods (difference-in-differences, regression discontinuity, instrumental variables) appear more credible not because they are assumption-free, but because their assumptions---parallel trends, continuity at the threshold, exclusion restrictions---are often more transparent and partially testable.
Selection-on-observables makes a stronger, less testable claim: That we have measured everything that matters.
But ``stronger'' does not mean ``never justified.''

\paragraph{A Pragmatic Perspective.}
If we insisted on ironclad causal identification for every empirical claim, much of applied research would grind to a halt.
Medicine advanced for centuries on observational evidence before randomized trials became standard.
Economics developed theories of market behavior, labor supply, and human capital from correlational data long before the credibility revolution.
Software engineering cannot always wait for perfect causally valid designs---practitioners need actionable insights now, even if those insights come with uncertainty.

The honest position is not that regression-based causal claims are invalid, but that they occupy a different place on the credibility spectrum.
A well-theorized regression study---one that explicitly articulates its causal model, defends the completeness of its adjustment set, addresses obvious threats, and conducts sensitivity analysis (e.g., Rosenbaum bounds~\cite{rosenbaum2002observational}, $E$-values~\cite{vanderweele2017sensitivity}, or coefficient stability tests~\cite{oster2019unobservable})---provides weaker but still informative evidence.
The key is transparency: Rather than claiming an effect is ``proven,'' the researcher specifies the conditions under which the estimate would be unbiased and invites scrutiny of whether those conditions hold.
Readers can then calibrate their confidence accordingly.
The danger lies not in using regression for causal questions, but in doing so naively.
The goal of this thesis is not to dismiss all prior work based on multivariate regressions, but to articulate clearly the assumptions underlying different methods, so that researchers can choose designs appropriate to their questions and honestly communicate the strength of their evidence.

\section{Formalizing Causal Inference}
\label{sec:formalism}

Having motivated the need for rigorous causal reasoning through the pitfalls of naive approaches, we now turn to the formal foundations.
The goal of formalization is not mathematical abstraction for its own sake, but precision: To state exactly what we mean by a ``causal effect,'' to articulate the assumptions required to estimate it, and to derive the logical consequences of those assumptions.
Without this precision, debates about causation devolve into disputes over intuition; with it, we can identify exactly where disagreements lie and what evidence would resolve them.

It is important to distinguish between \emph{frameworks} and \emph{methods}.
The frameworks presented in this section---potential outcomes and graphical causal models---are conceptual languages for reasoning about causation.
They define causal effects, formalize the assumptions required to identify them, and provide tools to assess their credibility.
However, they are not themselves statistical procedures that one applies to data.
They should be used to motivate the choice of specific \emph{methods}---randomized controlled trials, difference-in-differences, regression discontinuity, instrumental variables, panel regressions, etc.
Each method makes specific assumptions that can be understood within either framework; the frameworks provide the theoretical foundation that justifies the methods' operation when their assumptions hold.

Two complementary formal frameworks dominate modern causal inference.
The \emph{potential outcomes framework}, developed by Neyman~\cite{neyman1923application} and formalized by Rubin~\cite{rubin1974estimating}, defines causation through counterfactual comparison: What would have happened under alternative treatments?
The \emph{graphical causal model framework}, developed by Pearl~\cite{pearl2009causality}, represents causal relationships as directed graphs and provides algorithmic tools for determining when causal effects can be estimated from observational data.
Though they differ in their choice of formalism and emphasis, the two frameworks are mathematically compatible and often complementary.
In fact, applied researchers often use both to assess the causal credibility of their research designs (see Section~\ref{sec:pragmatic} on how to specifically apply them to the programming language vs. defect proneness example).

\subsection{The Potential Outcomes Framework}

The potential outcomes framework begins with a simple question: What does it mean to say that $X$ causes $Y$?
The answer, formalized through counterfactual reasoning, is that $X$ causes $Y$ for a particular unit if $Y$ would have been different had $X$ been different, holding all else constant.
This definition transforms causation from a vague philosophical concept into a precise comparison between two states of the world---one in which treatment occurs and one in which it does not.

\paragraph{Notation and Setup.}
Consider a population of units indexed by $i = 1, \ldots, N$.
Each unit is characterized by a treatment indicator $D_i \in \{0, 1\}$, where $D_i = 1$ indicates that unit $i$ received treatment and $D_i = 0$ indicates control.
The key conceptual move is to posit that each unit has \emph{two} potential outcomes:
\begin{itemize}[leftmargin=15pt]
    \item $Y_i(1)$: The outcome that \emph{would} obtain if unit $i$ receives treatment ($D_i = 1$)
    \item $Y_i(0)$: The outcome that \emph{would} obtain if unit $i$ does not receive treatment ($D_i = 0$)
\end{itemize}
These potential outcomes exist conceptually for every unit, regardless of which treatment the unit actually receives.
The observed outcome is determined by which potential outcome is ``revealed'' by the treatment assignment:
\begin{equation}
    Y_i = D_i \cdot Y_i(1) + (1 - D_i) \cdot Y_i(0)
\end{equation}
This switching equation makes explicit that we observe $Y_i(1)$ when $D_i = 1$ and $Y_i(0)$ when $D_i = 0$, but never both.
The individual causal effect for unit $i$ is defined as:
\begin{equation}
    \tau_i = Y_i(1) - Y_i(0)
\end{equation}
This quantity answers the question: How much would $i$'s outcome change if we switched $i$ from control to treatment?
The challenge is that $\tau_i$ is inherently unobservable.
We can observe at most one of $Y_i(1)$ or $Y_i(0)$ for any given unit; the other remains forever counterfactual.
A developer either uses one language or another to implement a specific program of interest.
We cannot rewind time and observe what would have happened under the alternative.

This impossibility of observing individual causal effects has profound implications.
It means that causal inference is fundamentally a problem of \emph{missing data}: For every treated unit, we are missing the control outcome; for every control unit, we are missing the treated outcome.
The entire enterprise of causal inference consists of finding ways to fill missing potential outcomes---or, more precisely, to estimate population summaries without needing to observe them individually.

\paragraph{Causal Estimands.}
Since individual effects are unobservable, we focus on population-level summaries.
The most common estimand is the \emph{Average Treatment Effect} (ATE):
\begin{equation}
    \text{ATE} = E[Y(1) - Y(0)] = E[Y(1)] - E[Y(0)]
\end{equation}
The ATE answers: On average across the population, how much does treatment change outcomes?
This is the quantity of interest when we ask whether a policy should be adopted universally.

An alternative estimand is the \emph{Average Treatment Effect on the Treated} (ATT):
\begin{equation}
    \text{ATT} = E[Y(1) - Y(0) \mid D = 1]
\end{equation}
The ATT answers: Among those who actually received treatment, what was the average effect?
This is relevant when evaluating an existing program: Given that certain people chose to participate, did participation help them?
Note that ATE and ATT coincide only when treatment effects are homogeneous or when treatment assignment is independent of potential outcomes.

Other estimands include the \emph{Average Treatment Effect on the Controls} (ATC), \emph{quantile treatment effects} (effects on specific percentiles of the outcome distribution), and \emph{conditional average treatment effects} (CATEs) that allow effects to vary with observed characteristics.
%The choice of estimand should be guided by the substantive question, not statistical convenience.

\paragraph{From Estimands to Identification.}
Having defined the estimands of interest, the next question is under what conditions they can be estimated from observed data.
Consider the naive estimator that compares mean outcomes between treated and control groups:
\begin{equation}
    E[Y \mid D=1] - E[Y \mid D=0]
\end{equation}
Expanding in terms of potential outcomes reveals:
\begin{align}
    E[Y \mid D=1] - E[Y \mid D=0] &= E[Y(1) \mid D=1] - E[Y(0) \mid D=0] \nonumber\\
    &= \underbrace{E[Y(1) \mid D=1] - E[Y(0) \mid D=1]}_{\text{ATT}} \\&+ \underbrace{E[Y(0) \mid D=1] - E[Y(0) \mid D=0]}_{\text{Selection Bias}}
    \label{eq:selection-bias}
\end{align}
The first term is the ATT.
The second is \emph{selection bias}: the difference in baseline potential outcomes between treated and untreated individuals.
When treatment is self-selected, these groups may differ systematically even in the absence of a treatment effect — high-ability developers may adopt functional programming languages precisely because they are high-ability, so comparing their outcomes with non-adopters conflates the language's effect with pre-existing ability differences.
Whether the naive comparison recovers a causal effect depends entirely on whether selection bias vanishes.
We now apply this reasoning to two canonical settings.

\paragraph{Case 1: Randomized Control Trials.}
In a randomized control trial (RCT), treatment is assigned by a mechanism that is independent of the units' characteristics, ensuring:
\begin{equation}
    (Y(1), Y(0)) \perp D
    \label{eq:rct-independence}
\end{equation}
Independence implies $E[Y(0) \mid D=1] = E[Y(0) \mid D=0] = E[Y(0)]$, so selection bias in Equation~\eqref{eq:selection-bias} is exactly zero.
Similarly, $E[Y(1) \mid D=1] = E[Y(1)]$, so the ATT equals the ATE, and the naive difference in means identifies the ATE:
\begin{equation}
    E[Y \mid D=1] - E[Y \mid D=0] = E[Y(1)] - E[Y(0)] = \text{ATE}
\end{equation}
Randomization alone suffices; no covariates, modeling assumptions, or functional-form restrictions are required.
This is why RCTs are considered the gold standard for causal inference: Identification follows from the study design rather than from the plausibility of assumptions.

A valid RCT additionally requires the \emph{Stable Unit Treatment Value Assumption (SUTVA)}---that unit $i$'s potential outcomes depend only on $i$'s own treatment status (no interference) and that treatment is well-defined (no hidden variations)---as well as \emph{compliance}: units must receive the treatment they are assigned.
If subjects opt out or cross over, the realized treatment diverges from the random assignment, and additional assumptions (e.g., monotonicity in the instrumental variables sense) are needed to recover causal effects.
Under randomization, SUTVA, and compliance, the sample difference $\bar{Y}_{\text{treated}} - \bar{Y}_{\text{control}}$ is an unbiased estimator of the ATE.

\paragraph{Case 2: Multivariate Regression.}
When randomization is infeasible---as is common in software engineering, where programming languages, tools, and team structures cannot be randomly assigned---researchers often turn to multivariate regression:
\begin{equation}
    Y_i = \alpha + \tau D_i + X_i' \beta + \varepsilon_i
    \label{eq:regression}
\end{equation}
where $X_i$ is a vector of observed covariates and $\varepsilon_i$ is the error term.
The potential outcomes framework makes explicit the assumptions under which the OLS estimate $\hat{\tau}$ admits a causal interpretation.
Specifically, the key assumption is \emph{conditional ignorability} (also called unconfoundedness or selection on observables):
\begin{equation}
    (Y(1), Y(0)) \perp D \mid X
    \label{eq:cond-ignorability}
\end{equation}
This asserts that, among units with the same observed characteristics $X$, treatment assignment is independent of potential outcomes---that is, treatment is as-if randomly assigned within strata defined by $X$.
A complementary assumption is \emph{overlap} (positivity): $0 < P(D=1 \mid X=x) < 1$ for all $x$ in the support of $X$, ensuring both treated and untreated units exist at every covariate profile.
Under these two conditions, the RCT logic applies within each stratum.
Selection bias vanishes conditionally---$E[Y(0) \mid D=1, X] = E[Y(0) \mid D=0, X]$---so:
\begin{equation}
    E[Y \mid D=1, X] - E[Y \mid D=0, X] = E[Y(1) - Y(0) \mid X]
\end{equation}
Averaging over $X$ recovers the ATE: $\text{ATE} = E_X\!\big[E[Y \mid D\!=\!1, X] - E[Y \mid D\!=\!0, X]\big]$.
The regression in Equation~\eqref{eq:regression} implements this identification result by modeling $E[Y \mid D, X]$ as linear and additive, which introduces a further assumption: that the functional form is correctly specified and that the treatment effect $\tau$ is constant across units.

In summary, for $\hat{\tau}$ to be causally interpretable, one must assume:
\begin{enumerate}[leftmargin=15pt]
    \item \textbf{SUTVA}: No interference between units and a well-defined treatment.
    \item \textbf{Conditional ignorability}: All variables that jointly affect treatment and outcome are included in $X$; no unmeasured confounders remain.
    \item \textbf{Overlap}: Both treatment groups are represented at all covariate values.
    \item \textbf{Correct functional form}: The linear, additive specification is a reasonable approximation to the true conditional expectation function.
\end{enumerate}
Unlike randomization, conditional ignorability is fundamentally untestable: we cannot verify from data alone that all confounders have been measured, since the confounders we might be missing are, by definition, unobserved.

%The potential outcomes framework makes these requirements explicit, forcing the analyst to articulate \emph{why} a particular regression warrants a causal interpretation rather than treating covariate adjustment as sufficient by itself.

\begin{comment}
\paragraph{Why Simple Comparisons Fail.}
A naive approach to estimating the ATE would compare mean outcomes between treated and control groups:
\begin{equation}
    E[Y \mid D=1] - E[Y \mid D=0]
\end{equation}
Expanding this in terms of potential outcomes reveals the problem:
\begin{align}
    E[Y \mid D=1] - E[Y \mid D=0] &= E[Y(1) \mid D=1] - E[Y(0) \mid D=0] \\
    &= \underbrace{E[Y(1) \mid D=1] - E[Y(0) \mid D=1]}_{\text{ATT}} \\ &+ \underbrace{E[Y(0) \mid D=1] - E[Y(0) \mid D=0]}_{\text{Selection Bias}}
\end{align}
The first term is the ATT---the causal effect among the treated.
The second term is \emph{selection bias}: The difference in baseline potential outcomes between those who select into treatment and those who do not.
If treatment is self-selected, these groups may differ systematically even in the absence of treatment.
High-ability developers may adopt new tools precisely because they are high-ability; comparing their outcomes with those of non-adopters confounds the tool's effect with pre-existing ability differences.
The observed difference equals the causal effect only when selection bias is zero---that is, when $E[Y(0) \mid D=1] = E[Y(0) \mid D=0]$.
This holds when treatment is randomly assigned, but rarely otherwise.

\paragraph{Identifying Assumptions.}
To bridge the gap between observable quantities and causal estimands, we require assumptions that discipline how potential outcomes relate to treatment assignment.
Three assumptions are foundational:

\begin{enumerate}[leftmargin=15pt]
    \item \textbf{Stable Unit Treatment Value Assumption (SUTVA).}
    This assumption has two components.
    First, \emph{no interference}: The potential outcomes for unit $i$ depend only on $i$'s own treatment status, not on the treatment status of other units.
    If one developer's adoption of a tool changes another developer's productivity (through collaboration, knowledge sharing, or competition), SUTVA is violated.
    Second, \emph{no hidden variations of treatment}: There is only one version of treatment, so $Y_i(1)$ is well-defined.
    If ``adopting a tool'' means different things for different developers (some use it intensively, others superficially), the single potential outcome $Y_i(1)$ may be poorly defined.
    SUTVA is often violated in networked settings.
    If developers on the same team influence each other, individual-level analysis may be inappropriate; the relevant unit may be the team, not the developer.
    
    \item \textbf{Ignorability (Unconfoundedness, Conditional Independence).}
    Conditional on observed covariates $X$, treatment assignment is independent of potential outcomes:
    \begin{equation}
        (Y(1), Y(0)) \perp D \mid X
    \end{equation}
    This assumption asserts that, among units with identical observed characteristics, treatment is as-if randomly assigned.
    Any systematic differences between treated and control groups are fully captured by $X$; once we condition on $X$, no unmeasured confounders remain.
    Ignorability is the critical assumption for observational studies.
    It is inherently untestable---because we cannot observe the unmeasured confounders, we cannot verify that all confounders have been measured.
    Its plausibility rests entirely on the researcher's knowledge of the data-generating process---or the plausibility of the causal theory based on the researcher's knowledge.
    
    \item \textbf{Overlap (Positivity, Common Support).}
    For all values of $X$ in the population, both treatment states have positive probability:
    \begin{equation}
        0 < P(D=1 \mid X) < 1 \quad \text{for all } X
    \end{equation}
    This ensures that for any combination of covariates, we can find both treated and untreated units, enabling comparison.
    If some covariate combinations are treated with probability 1 or 0, we have no counterfactual basis for estimating effects in that region.
    Overlap is testable: We can usually examine whether the distributions of covariates (or propensity scores) in treated and control groups have common support.
    Violations suggest that certain treatment effects are not identified from the data.
\end{enumerate}

\paragraph{Identification Under Ignorability.}
Under SUTVA, ignorability, and overlap, the ATE is identified from observed data.
The key insight is that ignorability allows us to substitute observable conditional expectations for unobservable potential outcome means:
\begin{align}
    E[Y(1)] &= E_X[E[Y(1) \mid X]] = E_X[E[Y(1) \mid X, D=1]] = E_X[E[Y \mid X, D=1]] \\
    E[Y(0)] &= E_X[E[Y(0) \mid X]] = E_X[E[Y(0) \mid X, D=0]] = E_X[E[Y \mid X, D=0]]
\end{align}
The first equality uses iterated expectations.
The second uses ignorability: $E[Y(1) \mid X] = E[Y(1) \mid X, D=1]$ because treatment is independent of potential outcomes given $X$.
The third uses the switching equation: Among the treated, observed $Y$ equals $Y(1)$.

Therefore:
\begin{equation}
    \text{ATE} = E_X[E[Y \mid D=1, X] - E[Y \mid D=0, X]]
\end{equation}
This expression involves only observable quantities: Conditional means of $Y$ given treatment status and covariates, averaged over the covariate distribution.
Various estimation strategies---regression adjustment, propensity score matching, inverse probability weighting, doubly robust estimation---implement this identification result using different statistical approaches.
\end{comment}

\subsection{Graphical Causal Models}

While the potential outcomes framework excels at defining estimands and clarifying the role of randomization, it is less helpful for reasoning about \emph{which} variables to condition on.
The graphical causal model framework, developed by Pearl and colleagues, addresses this gap by representing causal assumptions as directed acyclic graphs (DAGs) and providing algorithmic criteria for determining when causal effects are identifiable~\cite{pearl2009causality}.

\paragraph{Directed Acyclic Graphs.}
A DAG $G = (V, E)$ consists of a set of nodes $V$ representing variables and a set of directed edges $E$ representing direct causal relationships.
An edge $X \to Y$ asserts that $X$ has a direct causal effect on $Y$---changing $X$ while holding all other causes of $Y$ fixed would change $Y$.
The graph is ``acyclic'' because no variable can be its own ancestor: Following directed edges, one can never return to the starting node.
This rules out simultaneous causation (which requires more complex frameworks) but permits chains of causation over time.

The power of DAGs lies in what they make explicit.
Every arrow represents a causal claim; every absent arrow asserts the absence of a direct effect.
Critics can challenge a causal model by questioning specific edges: ``Why do you assume $X$ affects $Y$ directly?'' or ``Have you considered that $Z$ might cause both $X$ and $Y$?''
This transparency is impossible with regression equations, where causal assumptions are buried in the choice of variables and functional forms.

\paragraph{Paths and Associations.}
In a DAG, variables can be associated through multiple mechanisms.
A \emph{path} is any sequence of edges connecting two variables, regardless of direction.
Three types of paths are fundamental:

\begin{enumerate}[leftmargin=15pt]
    \item \emph{Causal paths} (or directed paths): $X \to \cdots \to Y$.
    All edges point from $X$ toward $Y$.
    These represent the causal effect of $X$ on $Y$, possibly mediated through intermediate variables.
    
    \item \emph{Back-door paths} (or confounding paths): $X \leftarrow \cdots \to Y$ or $X \leftarrow \cdots \leftarrow Z \to \cdots \to Y$.
    These paths have an arrow \emph{into} $X$ and connect $X$ to $Y$ through a common cause.
    They represent spurious association due to confounding.
    
    \item \emph{Collider paths}: $X \to \cdots \to C \leftarrow \cdots \leftarrow Y$.
    These paths pass through a \emph{collider}---a variable with arrows pointing into it from both directions.
    Collider paths are naturally ``blocked'' and do not transmit association unless we condition on the collider (or its descendants), at which point they become ``open.''
\end{enumerate}

\paragraph{d-Separation: When Variables Are Independent.}
The criterion of \emph{d-separation} (``d'' for directional) determines when two variables are statistically independent given a conditioning set.
A path between $X$ and $Y$ is \emph{blocked} by a set $Z$ if either: (1) The path contains a non-collider (a variable with at most one arrow pointing in) that is in $Z$, or (2) The path contains a collider that is \emph{not} in $Z$ and has no descendants in $Z$.
Variables $X$ and $Y$ are d-separated by $Z$ if all paths between them are blocked by $Z$.
D-separation implies conditional independence: If $X$ and $Y$ are d-separated by $Z$ in the graph, then $X \perp Y \mid Z$ in any distribution compatible with the graph.

The practical implication is that we can read off conditional independence relations directly from the graph, without computing probabilities.
This allows researchers to check whether their assumed causal structure is consistent with observed data (by testing the implied independencies) and to identify which variables must be conditioned on to block confounding.

\paragraph{The do-Operator and Interventional Distributions.}
The central distinction in Pearl's framework is between \emph{observing} that $X = x$ and \emph{intervening} to set $X = x$.
These correspond to different probability distributions:
\begin{itemize}[leftmargin=15pt]
    \item $P(Y \mid X = x)$: The conditional distribution of $Y$ given that we \emph{observe} $X = x$.
    This may reflect both the causal effect of $X$ on $Y$ and confounding from common causes.
    
    \item $P(Y \mid do(X = x))$: The distribution of $Y$ when we \emph{intervene} to set $X = x$, overriding whatever process normally determines $X$.
    This isolates the causal effect, as the intervention breaks the connection between $X$ and its causes.
\end{itemize}

Graphically, the intervention $do(X = x)$ corresponds to \emph{graph surgery}: We delete all arrows pointing into $X$ (since the intervention replaces $X$'s natural causes) while preserving all arrows pointing out of $X$ (since $X$ still affects its descendants).
In the modified graph, $X$ becomes an exogenous variable set to $x$, and the resulting distribution is the interventional distribution.

The fundamental goal of causal inference is to estimate $P(Y \mid do(X = x))$ from observational data, which only provides $P(Y \mid X = x)$.
When are these equal?
Precisely when there is no confounding---that is, when all paths from $X$ to $Y$ except the causal paths are blocked.

\paragraph{The Back-Door Criterion.}
Pearl's back-door criterion provides a graphical test for when conditioning on a set of variables $Z$ suffices to identify the causal effect.
A set $Z$ satisfies the back-door criterion relative to $(X, Y)$ if:
\begin{enumerate}[leftmargin=15pt]
    \item No variable in $Z$ is a descendant of $X$ (to avoid conditioning on mediators or colliders on the causal path).
    \item $Z$ blocks all back-door paths from $X$ to $Y$ (paths with an arrow into $X$).
\end{enumerate}

When the back-door criterion is satisfied, the interventional distribution is identified by adjustment:
\begin{equation}
    P(Y \mid do(X = x)) = \sum_z P(Y \mid X = x, Z = z) P(Z = z)
\end{equation}
This formula has a natural interpretation: We compute the effect of $X$ on $Y$ within each stratum defined by $Z$, then average over the distribution of $Z$ in the population.
The result is the causal effect, purged of confounding.

The back-door criterion illuminates the ``bad controls'' problem discussed earlier.
Conditioning on a mediator (a descendant of $X$ on the causal path) violates condition 1 and blocks part of the causal effect.
Conditioning on a collider opens a previously blocked path and introduces bias.
The graphical framework provides a systematic way to determine which variables should and should not be included in an adjustment set.

\paragraph{Beyond the Back-Door Criterion.}
When no set of observed variables satisfies the back-door criterion---perhaps because key confounders are unmeasured---causal effects may still be identifiable through alternative strategies.
Pearl's do-calculus provides a complete set of rules for determining identifiability from any combination of observational and experimental data.
The \emph{front-door criterion}, for instance, identifies causal effects through mediating variables when the direct path is confounded but the mediator-outcome relationship is not.
Instrumental variables, discussed in Section~\ref{sec:identification}, can be understood graphically as exploiting exogenous variation that affects treatment but not outcomes directly.

\subsection{Using the Frameworks in Practice}

How should an applied researcher use these frameworks?
In practice, neither framework is a statistical method that one ``runs'' on data.
Rather, they are \emph{conceptual tools for reasoning about causation}---languages for articulating what causal effect one seeks to estimate, what assumptions would justify that estimate, and what could go wrong.
The actual estimation is performed by specific statistical methods---difference-in-differences, regression discontinuity, instrumental variables, matching---which the following sections present.
Each method makes specific assumptions (parallel trends, continuity at the threshold, exclusion restrictions) that can be understood and justified within either framework.

The practical workflow typically proceeds as follows.
First, the researcher uses the potential outcomes framework to clarify the research question: What is the treatment?
What is the outcome?
What causal estimand (ATE, ATT, LATE) is of interest?
Second, the researcher draws a causal graph encoding domain knowledge about the data-generating process, identifying potential confounders, mediators, and colliders.
This step forces explicit articulation of assumptions that might otherwise remain implicit.
Third, the researcher selects an identification strategy---a research design that, under stated assumptions, provides a credible path from observed data to causal effects.
The choice depends on what natural experiments or exogenous variation the setting provides.
Finally, the researcher implements the chosen method using appropriate statistical techniques, conducts robustness checks, and interprets results in light of the assumptions made.

Though developed independently, the two frameworks are mathematically compatible and often used together.
The potential outcomes condition of ignorability---$(Y(1), Y(0)) \perp D \mid X$---corresponds precisely to d-separation in an appropriately constructed graph.
In practice, potential outcomes excel at the first step (defining estimands and clarifying what we want to learn), while graphical models excel at the second step (encoding assumptions visually and determining what to adjust for).
The frameworks thus serve as the \emph{theoretical foundation} upon which specific methods rest: A difference-in-differences analysis implicitly assumes that the parallel trends assumption blocks all back-door paths; an instrumental variable strategy assumes the instrument satisfies exclusion restrictions that can be depicted graphically.
Understanding both frameworks enables researchers to see what each method achieves and where it might fail---knowledge essential for choosing appropriate designs and interpreting results honestly.

\section{Threats to Causal Validity and Its Mitigations}
\label{sec:threats}

Every causal design, no matter how carefully constructed, faces threats that could undermine its conclusions.
The literature on research methodology distinguishes three fundamental types of validity: internal validity, which concerns whether the estimated effect reflects the true causal relationship within the study sample; external validity, which concerns whether findings generalize beyond the specific study context; and construct validity, which concerns whether the measured variables actually capture the theoretical constructs of interest.
Understanding these threats is essential both for designing rigorous studies and for critically evaluating the research of others.

\subsection{Internal Validity}

A study possesses high internal validity when we can confidently attribute observed outcome differences to the treatment rather than to confounding factors.
Again, the fundamental challenge is that we never observe the counterfactual and must instead infer it from comparison groups under assumptions that may or may not hold.

The most pervasive threat to internal validity is \emph{confounding}, which occurs when an unmeasured variable affects both treatment assignment and outcomes, creating a spurious association that masquerades as a causal effect.
Consider a study examining whether adopting a new development tool improves code quality.
If repositories with more experienced maintainers are both more likely to adopt innovative tools and more likely to produce high-quality code regardless of the tool, the observed association between tool adoption and quality may reflect maintainer expertise rather than the tool's genuine effect.
Randomized experiments eliminate confounding by construction, since random assignment ensures that treatment and control groups are, in expectation, identical on all characteristics---both observed and unobserved.
Observational designs lack this guarantee and must either measure and adjust for potential confounders or rely on design-based strategies such as difference-in-differences, regression discontinuity, or instrumental variables that address confounding through alternative identifying assumptions.

Closely related to confounding is \emph{selection bias}, which arises when the process that determines who enters the study or receives treatment is systematically related to potential outcomes.
If only repositories already experiencing quality problems adopt a static analysis tool (hoping it will help), naive comparisons of adopters to non-adopters will conflate the tool's effect with the pre-existing quality differential.
Selection can also manifest as differential attrition: If struggling projects are more likely to cease activity and drop out of a longitudinal study, the surviving sample becomes increasingly unrepresentative, and estimates based on this sample may be severely biased.

Even with perfect identification of causal effects, \emph{measurement error} in treatments or outcomes can attenuate or distort estimates.
When treatment status is measured imprecisely---for example, inferring ``tool adoption'' from the presence of a configuration file that may not indicate actual sustained use---the resulting measurement error typically biases estimates toward zero, making true effects appear smaller than they are.
When outcomes are measured with systematic error correlated with treatment status, bias can operate in either direction and is harder to diagnose.

Finally, the Stable Unit Treatment Value Assumption (SUTVA) requires that one unit's treatment status does not affect another unit's outcomes---an assumption frequently violated in software engineering contexts.
Developers share knowledge across teams, projects share dependencies through package ecosystems, and practices diffuse through professional networks and organizations.
If adopting a tool in one repository indirectly improves code quality in related repositories through shared learning or code reuse, simple treatment-control comparisons will underestimate the tool's total effect by ignoring these spillovers.
Alternatively, if control units are ``contaminated'' by exposure to treatment effects (perhaps through documentation or community discussions), the estimated treatment effect will be biased toward zero.

\subsection{External Validity}

A study can achieve perfect internal validity---correctly identifying the causal effect within its specific sample---yet provide limited guidance for other populations, settings, or time periods.
External validity concerns this generalizability, and its threats are often underappreciated in fields that prioritize identification over representativeness.

The most fundamental external validity concern is \emph{population generalizability}: effects estimated on one population may not apply to others with different characteristics.
Studies conducted on open-source GitHub repositories may not generalize to proprietary enterprise development, where organizational constraints, incentive structures, and development practices differ substantially.
Effects observed among early adopters of a technology---who tend to be more technically sophisticated and intrinsically motivated---may not hold for the broader population.
Volunteer participants in controlled experiments may differ systematically from developers who would never consent to participate, limiting the applicability of experimental findings.

\emph{Temporal generalizability} presents particular challenges in rapidly evolving technological domains.
The productivity impact of early AI coding assistants, measured with 2023-era tools on 2023-era codebases, may bear little resemblance to the effects of more advanced tools deployed in the late 2020s with more capable models.
Best practices validated through empirical research in one technological era may become obsolete---or even counterproductive---as the underlying technologies evolve.
This suggests that empirical software engineering findings should be understood as historically situated rather than as timeless truths.

Beyond population and temporal concerns, causal effects often exhibit substantial \emph{context dependence}, varying with organizational culture, developer expertise, project characteristics, programming language, or other factors in ways that limit straightforward transferability.
A tool that accelerates novice developers may impose coordination overhead that slows experts; a practice that enhances productivity in small, co-located teams may create bottlenecks at enterprise scale.
Recognizing this heterogeneity, researchers increasingly report not just average treatment effects but also how effects vary across observable subgroups, since understanding who benefits and who does not is often more actionable than knowing only the population average.

\subsection{Construct Validity}

Beyond the question of whether we have correctly identified a causal effect (internal validity) and whether it generalizes (external validity), construct validity asks whether the variables we measure actually capture the theoretical constructs we care about.
This challenge is particularly acute in software engineering, where many concepts of interest---productivity, code quality, technical debt, developer experience---are multidimensional abstractions that resist direct measurement.

The operationalization of treatments shapes what effects we estimate.
``Adopting continuous integration'' might be operationalized as setting up a CI system, actually running automated tests on every commit, or promptly responding to test failures---each capturing different aspects of the underlying practice.
Similarly, ``code quality'' might be measured through defect density, code review feedback, static analysis warnings, or maintainability indices, each reflecting different quality dimensions and potentially yielding different conclusions about treatment effects.

Even well-defined metrics may suffer from measurement imprecision.
Commit counts, used as a productivity proxy, conflate substantive work with commit granularity conventions that vary across developers and projects.
Issue counts conflate bugs with feature requests, user-reported problems with automated detections, and genuine defects with documentation gaps.
Lines of code added confounds productive output with code bloat.
These measurement limitations do not invalidate empirical research, but they counsel humility about the precision of our conclusions and underscore the value of triangulating across multiple operationalizations.

\subsection{Sensitivity Analysis and Robustness}

Because the identifying assumptions underlying causal inference are fundamentally untestable---we cannot directly verify that parallel trends would have held or that an instrument satisfies the exclusion restriction---rigorous research practice demands sensitivity analysis that explores how conclusions would change if assumptions were violated.

For matching and propensity score studies, Rosenbaum bounds quantify how strong an unmeasured confounder would need to be to explain away the observed effect~\cite{rosenbaum2002observational}.
If results remain robust to confounders implausibly stronger than any known variable, confidence in the findings increases; if a moderately strong confounder could reverse conclusions, appropriate caution is warranted.
For difference-in-differences, researchers examine sensitivity to pre-trend specifications: Event study plots visualizing pre-treatment coefficients help assess whether the parallel trends assumption is plausible, and researchers increasingly test whether conclusions survive alternative trend extrapolations.

A powerful approach to robustness is employing multiple approaches with different identifying assumptions.
For example, when two-way fixed effects, Callaway-Sant'Anna, and imputation estimators all yield consistent treatment effect estimates, the finding is unlikely to be an artifact of any single method's assumptions.
Conversely, divergent results across estimators signal that conclusions may depend on methodological choices and warrant careful investigation.

Placebo tests provide another diagnostic tool, examining whether effects appear where they should not.
If treatment effects appear \emph{before} treatment occurs, this suggests confounding rather than causation.
If treatment affects outcomes it theoretically should not influence---such as pre-determined characteristics fixed before treatment assignment---something is amiss with the research design.
Passing placebo tests does not prove validity, but failing them is highly informative about potential threats.

Ultimately, the credibility of causal claims rests not on asserting that assumptions are satisfied---a claim that can never be definitively verified---but on transparently examining their plausibility, demonstrating robustness to reasonable alternatives, and acknowledging the limitations that still remain in the current research design.

\section{A Pragmatic Approach to Causal Inference and Causal Credibility Assessments in Empirical SE Research}
\label{sec:pragmatic}

The preceding sections have presented causal inference with appropriate rigor, emphasizing the assumptions underlying each method and the threats that can undermine conclusions.
This rigor is essential---but it can also be paralyzing.
If every observational study suffers from potential confounding, if every natural experiment has debatable assumptions, and if even randomized experiments face external validity concerns, can we learn anything useful?

This section argues that we can, provided we approach causal inference pragmatically.
The goal is not epistemological perfection but actionable evidence that improves upon intuition and correlational analysis.
Software engineering researchers must navigate trade-offs thoughtfully, choosing appropriate methods for their contexts and communicating uncertainty honestly.

\subsection{The Internal-External Validity Trade-off}

The most fundamental trade-off in empirical research is between internal validity (did we correctly identify the causal effect in our study?) and external validity (does the effect generalize to settings we care about?).
These two goals are often in tension, and no study achieves perfection on both.

Consider a randomized controlled trial evaluating an AI coding assistant.
If we recruit volunteer participants, run the study in a controlled environment, and measure carefully defined tasks, we may achieve excellent internal validity---we know the tool caused the observed productivity difference.
But the participants may not represent typical developers; the tasks may not reflect real work; the environment may not capture the pressures and interruptions of actual practice.
The clean causal estimate may not generalize to contexts that matter.

Conversely, a large-scale observational study of AI tool adoption across thousands of GitHub repositories captures real-world variation and naturally represents diverse contexts.
But without randomization, we cannot be certain that observed productivity differences reflect the tool's effect rather than selection---perhaps more capable developers adopted the tool first because they are also more tech-savvy.
The estimate may be externally valid (it reflects what happens in the real world), but internally questionable (we do not know if the tool caused it).

This trade-off is not a flaw to be fixed but a reality to be navigated.
The goal is not to find the one perfect study but to build knowledge through multiple studies with different strengths, whose convergent findings provide confidence that transcends any single study's limitations.

\subsection{A Hierarchy of Evidence}

Rather than treating causal inference as all-or-nothing, it is useful to think of a spectrum of evidence quality.
Different research questions warrant different standards, and different contexts permit different designs.
At the top of the hierarchy sit randomized experiments with representative samples.
These studies provide strong internal validity and external validity when samples are representative.
Below this are natural experiments and quasi-experimental designs: regression discontinuity, instrumental variables, and well-executed difference-in-differences studies.
These achieve internal validity under specific assumptions that may be partially testable (e.g., parallel pre-trends, no manipulation at thresholds).
External validity depends on how representative the setting is.
Next come observational studies with careful control strategies: matching, propensity scores, fixed effects, and theory-driven covariate adjustment.
These reduce confounding but cannot eliminate it; internal validity depends on untestable assumptions about unobserved confounders.
At the bottom are descriptive and correlational studies, which document patterns but make no causal claims.
These are valuable for exploration, hypothesis generation, and understanding prevalence, but they cannot answer ``what if'' questions about interventions.

\paragraph{What the Hierarchy Does Not Mean.}
The hierarchy does not imply that lower-tier studies are worthless or that only randomized experiments deserve publication.
A thoughtful observational study of real-world tool adoption may provide more useful guidance than a randomized experiment on a contrived task with unrepresentative participants.
The hierarchy is about the strength of \emph{causal} evidence, not overall scientific value.
The hierarchy is also not about methodological purity.
A DiD study with obviously violated parallel trends is worse than a careful descriptive study that does not overclaim.
Execution quality matters as much as design choice.

\subsection{Principles for Pragmatic Practice}

Given these trade-offs, how should software engineering researchers approach causal inference in practice?
The answer lies not in a rigid hierarchy of methods but in thoughtful matching of research designs to research questions, contexts, and available resources.

The first principle is to \emph{match the design to the question and context}.
Not every research question requires---or can support---the strongest possible causal design.
For exploratory research aimed at generating hypotheses and identifying promising directions, descriptive and correlational studies remain entirely appropriate.
For consequential decisions with significant resource implications (e.g., Should an organization mandate a particular tool? Should a platform change its policies?), stronger causal evidence is warranted to justify the costs and risks of action.
The practical question is: Which designs are feasible given the specific setting? Is randomization possible, or would it be unethical or impractical? Does the context provide a natural experiment through policy discontinuities or staggered adoption? What data are available, and what quality are they?
Researchers should choose the strongest design that remains credible in their specific context, rather than forcing an inappropriate method that offers theoretically stronger causal identification but rests on implausible assumptions in the specific study setting.

The second principle is \emph{transparency about assumptions and limitations}.
Every causal claim, whether from an RCT or an observational study, rests on assumptions that cannot be definitively verified.
Rigorous research practice demands that these assumptions be stated explicitly, their plausibility discussed in the specific research context, and potential violations acknowledged forthrightly.
Transparency about limitations does not weaken a research contribution; rather, it strengthens it by demonstrating that the authors understand the nature of the evidence they are providing and the conditions under which their conclusions hold.

Third, because identifying assumptions are fundamentally untestable, researchers should \emph{conduct systematic sensitivity analysis} exploring how conclusions would change if assumptions were violated.
For matching studies, Rosenbaum bounds quantify the strength of unmeasured confounding required to overturn findings.
For difference-in-differences, researchers should examine robustness to alternative pre-trend specifications and conduct event study analyses.
Applying multiple estimators with different identifying assumptions---and reporting whether they yield consistent conclusions---provides additional evidence about the robustness of findings.
When conclusions survive these diagnostic checks, confidence in their validity grows; when they prove fragile, appropriate caution and qualification are warranted.

Fourth, researchers should \emph{embrace triangulation} as a strategy for building cumulative knowledge.
No single study, however well-designed and carefully executed, provides definitive evidence about causal relationships.
Scientific knowledge accumulates through replication across different contexts, methods, and research teams.
When randomized experiments, quasi-experimental observational studies, and qualitative research all point in the same direction, our collective confidence in the underlying causal relationship grows substantially.
When different approaches yield divergent conclusions, the discrepancy itself is informative: perhaps effects are genuinely context-dependent, varying across populations or settings in ways that merit investigation; or perhaps some studies contain unrecognized biases that others avoid.
Either way, the apparent contradiction advances understanding more than any single study could.

Fifth, researchers should \emph{report effect heterogeneity} rather than focusing exclusively on average treatment effects.
Population averages often conceal substantial variation across subgroups that is more actionable than the overall mean.
A tool that improves productivity by 10\% on average might help novices by 30\% while actually slowing experienced developers by 10\%---information far more useful to practitioners than the aggregate figure alone.
Examining how effects vary across observable characteristics---developer experience, project type, programming language, organizational context---provides insights that help readers assess whether findings are likely to generalize to their own situations.

Finally, researchers should \emph{communicate uncertainty honestly}.
Statistical significance is neither equivalent to practical importance nor a guarantee of correctness.
Reporting confidence intervals conveys the precision of estimates; discussing effect magnitudes in substantively meaningful terms helps readers evaluate practical relevance; acknowledging when evidence is suggestive rather than definitive maintains scientific integrity.
Overclaiming---presenting tentative findings as established facts---undermines the credibility of individual researchers and the field as a whole.
Appropriate epistemic humility, by contrast, builds trust and enables research to inform practice even when certainty remains elusive.

\subsection{The Value of Imperfect Evidence}

Given the stringent assumptions underlying causal inference and the many threats to validity that observational studies face, it might be tempting to conclude that without a perfect randomized experiment, we cannot learn anything meaningful about causation.
This conclusion, while superficially rigorous, is ultimately too pessimistic to serve as a guide for research practice.
Imperfect evidence, accumulated thoughtfully across multiple studies and methods, can guide decisions far better than no evidence at all or purely intuitive reasoning based on anecdote and conventional wisdom.

Consider a practical decision: Should a software organization adopt a new static analysis tool?
A correlational study finds that repositories currently using the tool exhibit 20\% fewer bugs than those that do not.
This association, while suggestive, is not causal evidence---perhaps better-organized teams both adopt sophisticated tooling and write higher-quality code, with no genuine effect of the tool itself.
But suppose a difference-in-differences study exploiting staggered adoption finds that bug rates drop by 15\% after repositories adopt the tool, with no evidence of differential pre-trends.
A propensity score matching study comparing adopters to similar non-adopters finds a 12\% reduction.
Qualitative interviews with developers reveal that the tool catches common coding mistakes that previously slipped through review.

No single study in this portfolio is methodologically definitive.
The DiD study might suffer from time-varying confounders; the matching study cannot account for unobserved selection; the interviews provide mechanism but not magnitude.
Yet together, these studies suggest that the tool probably provides genuine benefits, though the effect size might be smaller than naive correlations suggest and likely varies across contexts and development practices.
This is useful guidance for organizational decision-making---not certainty, but evidence that substantially improves upon intuition alone.

The goal of causal inference in software engineering is not to achieve epistemological perfection---an unattainable standard in any empirical science---but to provide better evidence than we would have otherwise.
By employing the most rigorous designs feasible in each context, maintaining transparency about assumptions and limitations, conducting systematic sensitivity analyses, and triangulating across complementary studies, researchers can contribute genuine knowledge to a field where randomized experiments remain rare and observational data are inevitably imperfect.
The credibility of causal claims ultimately rests not on claiming that assumptions are satisfied, but on transparently examining their plausibility, demonstrating robustness to reasonable violations, and honestly communicating what the evidence can---and cannot---tell us about the phenomena under investigation.
